\section{Holography}

Holography is the statement that everything inside a region is encoded on its boundary, and on this mesh it is not a
deep mystery but a plain consequence of the exponential growth. The bulk swells faster than its surface, so almost every
cell in a region lies near the edge, and information about the interior is crowded outward onto the boundary with
nowhere else to go. This gives three measured laws.

The first law is the area law. The ground state of the emergent field on the mesh has its entanglement entropy
growing with the AREA of a region's boundary, not its volume, while a hot, disordered state grows with the volume. This
is the signature of a local, low-energy state on a holographic geometry, and an experiment measures exactly this: it
computes the entanglement of a region in the ground state and finds it scales with the boundary area, then computes a
thermal state and finds it scales with the volume, the control that gives the area law its meaning.

The second law is the logarithmic boundary law. For a one-dimensional boundary interval, the entanglement across it
grows as the logarithm of the interval length, the Ryu-Takayanagi law, where the entangling surface is the shortest
path through the bulk between the interval's ends. An experiment measures the boundary-interval entanglement on the
mesh and finds the logarithm, while the same measurement on a flat control grows linearly instead, so the logarithm
is a property of the curved bulk, not of any boundary. The shortest bulk path between two boundary points is a genuine
geometric shortcut: two points far apart along the surface are only a few steps apart through the deep interior, which
an experiment also measures directly, joining distant boundary cells by a short hidden bulk path.

The third piece is the gravity propagator, the bulk-mediated correlation between two boundary points. Because the bulk
branches exponentially, locally like a tree, this correlation has an exact closed form: at the massless point of the
propagator it falls off as exactly $1/r^2$, fixed entirely by how fast the mesh branches. An experiment computes
this on the exact infinite tree by a cavity recursion and finds the $1/r^2$ law at both the four-dimensional mesh's
branching and its three-dimensional sibling's, with a massive field decaying faster, and validates the recursion
against a directly solved finite tree, where the level-to-level decay converges to the predicted value as the boundary
recedes. So holography on this mesh is three concrete, measured laws, the area law, the logarithmic entanglement, and
the $1/r^2$ boundary propagator, each a direct consequence of the exponential growth, and each demonstrated by an
experiment with a control where the answer comes out no.
